\documentclass[10pt]{article}
\usepackage[margin=0.85in]{geometry}
\usepackage{amsmath,amssymb}
\usepackage[hidelinks]{hyperref}

\begin{document}
\begin{center}
  {\Large\bfseries Even Euler characteristic and codimension-one immersion\par}
  \smallskip
  Disproof of Conjecture 00000003272
\end{center}

The closed, connected, oriented smooth six-manifold
\[
  M=\mathbb{CP}^{2}\times S^2
\]
has an even Euler characteristic and does not smoothly immerse
in \(\mathbb R^7\). Indeed, \(\mathbb{CP}^2\) has one cell in each
of dimensions \(0,2,4\), and \(S^2\) has one cell in each of
dimensions \(0,2\). Their product has six cells, all of even
dimension, so \(\chi(M)=6\).

\section*{The restricted tangent bundle}

Choose a projective line \(L=\mathbb{CP}^{1}\subset\mathbb{CP}^{2}\)
and a point \(p\in S^2\). Let \(j:L\longrightarrow M\) be
\(j(x)=(x,p)\). Write \(\varepsilon^r\) for the trivial real
rank-\(r\) bundle and \(H=\mathcal O(1)|_L\) for the dual
tautological complex line on \(L\). The product tangent bundle
restricts to
\[
  j^*TM\cong (T\mathbb{CP}^{2}|_L)_{\mathbb R}\oplus\varepsilon^2.
\]

To identify the right side, let \(\lambda\) be the tautological line
over \(\mathbb{CP}^2\). The complex tangent space at a line \(\ell\) is
\(\operatorname{Hom}_{\mathbb C}(\ell,\mathbb C^3/\ell)\), as seen by
writing nearby lines as graphs of linear maps into a complement.
Applying \(\operatorname{Hom}_{\mathbb C}(\lambda,-)\) to the
tautological inclusion gives the Euler sequence
\[
  0\longrightarrow \underline{\mathbb C}
   \longrightarrow (\lambda^*)^{\oplus3}
   \longrightarrow T\mathbb{CP}^2
   \longrightarrow 0.
\]
The first term consists of scalar endomorphisms of \(\lambda\).
This sequence splits as smooth complex bundles by choosing an
orthogonal complement to the first term. Restricting to \(L\)
and passing to real bundles gives
\[
  (T\mathbb{CP}^2|_L)_{\mathbb R}\oplus\varepsilon^2
  \cong H_{\mathbb R}^{\oplus3}.
\]

\section*{The obstruction}

The first Chern class of \(H\) generates \(H^2(L;\mathbb Z)\), and
its mod-two reduction \(\bar h\) generates \(H^2(L;\mathbb F_2)\).
For the underlying real bundle, \(w_2(H_{\mathbb R})=\bar h\)
\cite[Theorems~3.1--3.2 and Proposition~3.8]{hatcher}. Since
\(H^1(L;\mathbb F_2)=0\), all first Stiefel--Whitney classes
on \(L\) vanish. The degree-two Whitney sum formula now gives
\[
  w_2(j^*TM)=3\bar h=\bar h\ne0.
\]

An immersion \(M\longrightarrow\mathbb R^7\) would induce
\(TM\oplus\nu\cong\varepsilon^7\), where the normal bundle \(\nu\)
is a real line bundle. Pulling back along \(j\) would give
\(j^*TM\oplus j^*\nu\cong\varepsilon^7\) on \(L\). A line bundle
has zero second Stiefel--Whitney class, so the degree-two Whitney
formula on \(L\) would give
\[
  0=w_2(\varepsilon^7)=w_2(j^*TM)+w_2(j^*\nu)=\bar h,
\]
a contradiction. Thus even Euler characteristic is insufficient.

\section*{Lean verification}

The Lean~4.19.0 project computes the Euler characteristic from
the product cell dimensions. It models bundle isomorphism classes
on \(L\) with the standard degree-two characteristic class identities.
From the product restriction and the split Euler sequence it derives
the nonzero mod-two class and rules out a rank-one normal splitting.
The geometric identifications and characteristic class rules used
in this implication are established above. Run \verb|lake build| in
\verb|lean/| to verify the formal calculation.

\begin{thebibliography}{1}
\bibitem{hatcher}
A.~Hatcher, \emph{Vector Bundles and K-Theory}, version 2.2, 2017,
\href{https://pi.math.cornell.edu/~hatcher/VBKT/VB.pdf}{Section 3.1}.
\end{thebibliography}

\end{document}
